\documentclass[12pt]{article}
\usepackage[utf8]{inputenc}
\usepackage[T1]{fontenc}
\usepackage{hyperref}
\title{Teshenawa Language Resources for K-12 Curriculum Development}
\author{Compiled by Claude Sonnet 5 (claude-sonnet-5) for Andr\'as Kornai}
\date{}
\begin{document}
\maketitle

\section{Introduction}
Teshenawa (ISO 639-3 twc; Glottocode tesh1239) is an extinct Chadic (Afro-Asiatic) language of northern Nigeria. According to Wikipedia it was spoken in Teshena town, Kafin Hausa LGA, Jigawa State, and belongs to the Bade group of West Chadic. Its closest relative, Auyokawa, is also extinct. Wikipedia gives its extinction as after 1927. Searches found no living community of speakers and no teaching materials.

\section{Primary Educational Resources}
None found. No primers, textbooks or school programmes for Teshenawa turned up in search.

\section{Linguistic and Reference Materials}
\begin{itemize}
\item Teshenawa language (Wikipedia overview): \url{https://en.wikipedia.org/wiki/Teshenawa_language}
\item Bade languages (Wikipedia), which places Teshenawa within the Bade group: \url{https://en.wikipedia.org/wiki/Bade_languages}
\item Michael Bro\ss, 1997, a paper on the history of the extinct Auyo, Shira and Teshena languages (Glottolog bibliographic record only; I did not see the text): \url{https://glottolog.org/resource/reference/id/10548}
\item Michael Bro\ss, 1996, a paper on the classification of Shiranci and Teshenanci (Glottolog record only): \url{https://glottolog.org/resource/reference/id/124081}
\end{itemize}

\section{Religious and Community Resources}
None found.

\section{Digital and Audio Resources}
None found. No recordings or corpora were seen in search results.

\section{Curriculum Development Context}
Very little is documented about Teshenawa, and no speakers are known. A school curriculum cannot be built on the language itself from what was found. Realistic options are a short awareness unit on the Chadic languages of Jigawa State and on language extinction, drawing on the Wikipedia pages and the Bro\ss{} papers. Any further work would need archival research on the original 19th and early 20th century wordlists, which this search did not locate.
\end{document}